% Part: lambda-calculus
% Chapter: representability
% Section: currying

\documentclass[../../../include/open-logic-section]{subfiles}

\begin{document}

\olfileid{lam}{rep}{cur}
\olsection{కరీకరణ}

ఒక $\lambd$-అమూర్తీకరణ $\lambd[x][M]$ ఒకే ఆర్గ్యుమెంట్ గల
ప్రమేయాన్ని సూచిస్తుంది. అనేక ఆర్గ్యుమెంట్లను స్వీకరించే ప్రమేయాన్ని
నిర్వచించాలనుకున్నప్పుడు ఇది గణనీయమైన పరిమితి. దీన్ని చేయడానికి ఒక
మార్గం జతలు, త్రయాలు మొదలైనవాటిని నిర్మించడానికి వీలుగా
$\lambd$-కలనశాస్త్రాన్ని విస్తరించడం. అప్పుడు, ఉదాహరణకు, ఒక త్రిస్థానిక
ప్రమేయం $\lambd[x][M]$ తన ఆర్గ్యుమెంట్ ఒక త్రయం కావాలని ఆశిస్తుంది.
అయితే దీన్ని \emph{కరీకరణ} ద్వారా చేయడం మరింత అనుకూలం.

ఒక ఉదాహరణను పరిశీలిద్దాం. $\lambd$-కలనశాస్త్రంలో మనకు $+$ పరిక్రియ
ఉన్నట్లు కాసేపు ఊహిద్దాం. కూడిక ప్రమేయం $2$-స్థానికం, అంటే అది రెండు
ఆర్గ్యుమెంట్లను తీసుకుంటుంది. కానీ ఒక $\lambd$-అమూర్తీకరణ మనకు ఒకే
ఆర్గ్యుమెంట్ గల ప్రమేయాలను మాత్రమే ఇస్తుంది: $\lambd[(x,y)][(x+y)]$
వంటి వ్యక్తీకరణలను వాక్యనిర్మాణం అనుమతించదు. అయితే తన ఏకైక
ఆర్గ్యుమెంట్~$y$కు $x$ను కలిపే $\lambd[y][(x+y)]$ ద్వారా ఇచ్చిన
ఏకస్థానిక ప్రమేయం~$f_x(y)$ను పరిశీలించవచ్చు. నిజానికి ఇది ఒకే ప్రమేయం
కాదు; ప్రతి సంఖ్య~$x$కు ఒకటి చొప్పున ``$x$ను కలుపు'' అనే వేర్వేరు
ప్రమేయాల కుటుంబం. ఇప్పుడు $\lambd[x][f_x]$గా మరో ఏకస్థానిక
ప్రమేయం~$g$ను నిర్వచించవచ్చు. $x$ అనే ఆర్గ్యుమెంట్‌కు ప్రయోగిస్తే,
$g(x)$ ప్రమేయం~$f_x$ను ఫలితంగా ఇస్తుంది---అంటే దాని విలువలు ఇతర
ప్రమేయాలు. ఇప్పుడు $g$ను $x$కు, ఆ ఫలితాన్ని~$y$కు ప్రయోగిస్తే
$(g(x))y = f_x(y) = x+y$ వస్తుంది. ఈ విధంగా ఏకస్థానిక ప్రమేయం~$g$
ద్విస్థానిక కూడిక ప్రమేయం చేసే పనినే చేయగలదు. ``కరీకరణ'' అంటే
ద్విస్థానిక ప్రమేయాలను (వాటి విలువలే ఏకస్థానిక ప్రమేయాలుగా ఉండే)
ఏకస్థానిక ప్రమేయాలుగా మార్చే ఈ ఉపాయం.

ఇప్పుడు $\lambd$-కలనశాస్త్ర వాక్యనిర్మాణంలోనే ఒక ఉదాహరణ చూద్దాం.
$f(x,y) = x$ ప్రమేయాన్ని ఎలా సూచించాలి? రెండు ఆర్గ్యుమెంట్లను
స్వీకరించి మొదటిదాన్ని ఫలితంగా ఇచ్చే ప్రమేయాన్ని నిర్వచించాలంటే,
$\lambd[x][\lambd[y][x]]$ అని రాయవచ్చు. అక్షరార్థంగా ఇది ఒక
ఆర్గ్యుమెంట్~$x$ను స్వీకరించి, ప్రమేయం~$\lambd[y][x]$ను ఫలితంగా ఇచ్చే
ప్రమేయం. $\lambd[y][x]$ మరో ఆర్గ్యుమెంట్~$y$ను స్వీకరిస్తుంది; కానీ
దాన్ని విస్మరించి, ఎల్లప్పుడూ~$x$నే ఫలితంగా ఇస్తుంది. ఇప్పుడు
$\lambd[x][\lambd[y][x]]$ను రెండు ఆర్గ్యుమెంట్లకు ప్రయోగిస్తే ఏమవుతుందో
చూద్దాం:
\begin{align*}
  (\lambd[x][\lambd[y][x]])MN
  \bredone & (\lambd[y][M])N \\
  \bredone & M
\end{align*}

సాధారణంగా, ఏదో పదం~$N$ ద్వారా నిర్వచించిన $x_1$, \dots,~$x_n$
పరామితులు గల ప్రమేయాన్ని రాయడానికి
$\lambd[x_1][\lambd[x_2][\ldots\lambd[x_n][N]]]$ అని రాయవచ్చు. దానికి
$n$ ఆర్గ్యుమెంట్లను ప్రయోగిస్తే:
\begin{multline*}
  (\lambd[x_1][\lambd[x_2][\ldots\lambd[x_n][N]]]) M_1 \dots M_n \\
  \begin{aligned}
  \bredone {} & (\Subst{(\lambd[x_2][\ldots\lambd[x_n][N]])}{M_1}{x_1}) M_2
  \dots M_n\\
   \eqs {} & (\lambd[x_2][\ldots\lambd[x_n][\Subst{N}{M_1}{x_1}]]) M_2
            \dots M_n \\
  \vdots & \\
  \bredone {} &\Subst{\Subst{N}{M_1}{x_1}\ldots}{M_n}{x_n}
  \end{aligned}
\end{multline*}
\sourcecorrection{OLTELAMINT-001}{ఆంగ్ల మూలంలోని కరీకరణ తగ్గింపు శ్రేణి మొదటి పంక్తి చివర ఒక అదనపు $\bredone$ను పెట్టి, చివరి పంక్తిలో నిర్వచించిన పదం $N$కు బదులు నిర్వచించని $P$లో ప్రతిస్థాపించింది. శ్రేణిలో ఒక్క తగ్గింపు బాణమే ఉండేలా అదనపు బాణాన్ని తొలగించి, చివరి పంక్తిలో $N$ను పునరుద్ధరించాం.}
చివరి పంక్తి అక్షరార్థంగా ప్రమేయ నిర్వచనపు అంతర్భాగంలో $x_i$ స్థానంలో
$M_i$ని ప్రతిస్థాపించడాన్ని సూచిస్తుంది; అనేక ఆర్గ్యుమెంట్లను ప్రమేయానికి
ప్రయోగించినప్పుడు మనకు కావలసింది సరిగ్గా ఇదే.
\end{document}
